\fichatitulo{42 · FUNDAMENTO E RISCOS}{Relatório / preprint v1 · 2021}{Oportunidades e riscos dos modelos de base}
{\small Bommasani et al. \textit{On the Opportunities and Risks of Foundation Models}. Relatório / preprint v1 · 2021. \href{https://arxiv.org/pdf/2108.07258v1}{Fonte consultada}.\par}

\campo{Problema e objetivo}
Quais consequências decorrem de reutilizar modelos de base em muitas aplicações?

\campo{Principal contribuição · síntese interpretativa}
Articula oportunidades e riscos da emergência de capacidades e da homogeneização de componentes tecnológicos.

\campo{Método / natureza da evidência}
Relatório multidisciplinar; o recorte focal aborda conceito e dificuldades de avaliação de modelos adaptáveis.

\campo{Resultado ou argumento principal}
Distingue avaliação do modelo de base e avaliação de sistemas especializados, dependentes também do mecanismo de adaptação.

\campo{Excertos-chave · seleção literal}
\begin{itemize}
\excerto{1}{single points of failure}{Introdução.}
\excerto{2}{emergence and homogenization}{Resumo / introdução.}
\end{itemize}

\campo{Limitações e análise crítica}
Análise da ficha: o relatório não constitui um experimento único que quantifique todos os riscos. Reutilização ampla pode propagar falhas, mas o impacto precisa ser examinado em cada aplicação.

\campo{Aproveitamento na dissertação · proposta}
Fundamentação e discussão: avaliar o sistema RAG completo e documentar dependências de modelos compartilhados.

\registro{Fonte consultada nesta preparação: Resumo, seção 1 e abertura da seção 4.4 (avaliação); relatório não lido integralmente. Chave: \texttt{bommasaniEtAl2021foundationModels}.}
